% ─────────────────────────────────────────────────────────────
%  02-activite.tex — fiche de présentation d'une activité
%  mathématique pour le secondaire (2 à 3 pages).
%  À compiler avec pdfLaTeX (réglage par défaut d'Overleaf).
% ─────────────────────────────────────────────────────────────
\documentclass[11pt]{article}

\usepackage[T1]{fontenc}
\usepackage[french]{babel}
\usepackage{amsmath, amssymb}
\usepackage{siunitx}
\sisetup{output-decimal-marker = {,}}
\usepackage[margin=2.2cm]{geometry}
\usepackage{tikz}                 % figures dessinées en LaTeX
\usepackage{enumitem}             % réglage des listes
\usepackage{fancyhdr}             % en-tête et pied de page
\pagestyle{fancy}
\fancyhf{}
\lhead{Mathématique, 3\textsuperscript{e} secondaire}
\rhead{Activité : Pythagore au salon}
\cfoot{\thepage}

\title{Pythagore au salon : la taille d'un téléviseur}
\author{Prénom Nom}
\date{}

\begin{document}

\maketitle
\thispagestyle{fancy}

\section{Présentation de l'activité}

\begin{description}[leftmargin=3.6cm, style=nextline, font=\normalfont]
  \item[Niveau visé] 3\textsuperscript{e} secondaire (à adapter selon
    le groupe).
  \item[Durée] Une période de \qty{75}{\minute}.
  \item[Intention] Réinvestir la relation de Pythagore dans un contexte
    réel, et travailler les rapports et la proportionnalité.
  \item[Préalables] Relation de Pythagore, racine carrée, rapport et
    proportion.
  \item[Matériel] Calculatrice, ruban à mesurer, publicités de
    téléviseurs.
\end{description}

\section{Mise en situation}

Les téléviseurs sont vendus selon la longueur de leur diagonale,
exprimée en pouces ($\qty{1}{po} = \qty{2,54}{\centi\metre}$).
Un écran « 55 pouces » mesure donc \qty{55}{po} d'un coin à
l'autre. Mais quelle place occupera-t-il réellement sur le meuble du
salon ?

Les écrans actuels ont un format $16:9$ : la largeur et la hauteur
sont dans le rapport $16$ à $9$.

\begin{center}
\begin{tikzpicture}[scale=0.11]
  \draw[thick] (0,0) rectangle (48,27);
  \draw[dashed] (0,0) -- node[sloped, above, fill=white, inner sep=2pt]
    {diagonale \qty{55}{po}} (48,27);
  \node[below] at (24,0) {largeur $16k$};
  \node[right] at (48,13.5) {hauteur $9k$};
\end{tikzpicture}
\end{center}

\section{Déroulement}

\begin{enumerate}[label=\arabic*.]
  \item Mise en situation en grand groupe (\qty{10}{\minute}) :
    les élèves estiment la largeur d'un écran de \qty{55}{po}.
  \item Travail en équipe (\qty{35}{\minute}) : tâche de la section 4.
  \item Retour collectif (\qty{20}{\minute}) : comparaison des
    démarches, validation à l'aide d'un ruban à mesurer.
  \item Synthèse (\qty{10}{\minute}) : ce que la diagonale ne dit pas
    sur un écran.
\end{enumerate}

\section{Tâche de l'élève}

\begin{enumerate}[label=\alph*)]
  \item Montre que la diagonale d'un écran $16:9$ vaut
    $k\sqrt{337}$ si sa largeur est $16k$ et sa hauteur, $9k$.
  \item Calcule la largeur et la hauteur, en centimètres, d'un écran de
    \qty{55}{po}.
  \item Le meuble du salon mesure \qty{120}{\centi\metre} de large.
    L'écran de \qty{55}{po} y entre-t-il ? Justifie ta réponse.
  \item Un écran de \qty{110}{po} a une diagonale deux fois plus
    longue. Son aire est-elle deux fois plus grande ? Explique.
\end{enumerate}

\section{Corrigé et pistes d'intervention}

\begin{enumerate}[label=\alph*)]
  \item Par Pythagore,
    $d^2 = (16k)^2 + (9k)^2 = 256k^2 + 81k^2 = 337k^2$, donc
    $d = k\sqrt{337}$.
  \item $k = \dfrac{55}{\sqrt{337}} \approx \num{2,996}$, d'où une
    largeur de $16k \approx \qty{47,9}{po} \approx \qty{121,8}{\centi\metre}$
    et une hauteur de $9k \approx \qty{27,0}{po} \approx
    \qty{68,5}{\centi\metre}$.
  \item Non : il manque près de \qty{2}{\centi\metre}, sans compter le
    cadre de l'écran.
  \item Non : les deux dimensions doublent, donc l'aire est multipliée
    par $2^2 = 4$.
\end{enumerate}

\noindent Pistes d'intervention : un élève qui calcule $55 \div 16$
confond la diagonale avec la largeur ; lui faire tracer le rectangle à
l'échelle. La question d) prépare le lien entre rapport de similitude
et rapport des aires.

\end{document}
